\section*{September 15, 2026: PPML and dollar predictions from log prices}
\textit{Drafted by Codex.}

The two requested price-scale checks use the existing half-mile single-family
sample for 2008--2018. Annual regressions retain 11,202 sales at 68 sites;
pooled DiDs omit 2013 and retain 10,159. School-site and year fixed effects,
transaction weighting, property controls, and site-clustered intervals match
the preceding comparisons. All four original OLS event series are reproduced.

With property controls, the pooled PPML estimate is 2.06\% (95\% interval:
$-7.61\%$ to 12.74\%). The OLS-log coefficient transforms to 2.98\%
($-12.20\%$ to 20.77\%); dollar OLS estimates \$34,851
($-\$18,377$ to \$88,079). PPML does not reveal a strong positive effect
missing from logs. Its outcome is dollar price with an exponential conditional
mean; using PPML does not require integer prices or a Poisson outcome
distribution. Price is scaled to thousands for numerical stability.

\includegraphics[width=0.95\textwidth]{../input/price_scale_ppml.png}

\small PPML follows \href{https://personal.lse.ac.uk/tenreyro/lgw.html}{Santos Silva and Tenreyro (2006)}.
A difference between PPML and OLS logs would not alone identify a right-tail
composition mechanism. The estimators impose different conditional-mean
specifications and parallel-trend restrictions.
\normalsize
\clearpage

\subsection*{The log model can generate positive additive dollar estimates}
Dollar predictions multiply exponentiated fitted log prices by the pooled
mean of exponentiated log residuals (\href{https://doi.org/10.1080/01621459.1983.10478017}{Duan smearing}).
This assumes a common retransformation factor; heteroskedasticity can invalidate
that assumption. A sensitivity holds separate treated/control factors from
pre-2013 residuals fixed over time. No group-year factor is fitted to match
observed dollar means. All predictions are in-sample, with no uncertainty bands.

Regressing these predictions on the original levels specification yields a
2018 estimate of \$38,724 with property controls, compared with the observed
\$33,344. Fixed pre-period group smearing yields \$40,443. The main
reconstruction reduces squared error relative to zero event coefficients by
65\% over 2014--2018, but misses some annual moves. This is a descriptive
reconstruction metric, not the share of a causal effect explained. The raw
2018 mean gap is \$76,475 versus a predicted \$99,140.

Removing the log treatment-year terms while retaining all other fitted terms
produces an even larger 2018 dollar projection (\$104,145). The direct
contrast from keeping those terms is negative (about $-\$58,067$ among treated
sales in 2018). The positive additive projection therefore should not be read
as a positive treatment effect implied by the log model. Common proportional
movements can create different dollar movements from different baseline
prices. Observed sale mix and functional form also enter this reconstruction.
The pre-trend, anticipation, and selection concerns remain unresolved.

\includegraphics[width=0.92\textwidth]{../input/price_scale_predictions.png}

\small Reproduction: run \texttt{make} in
\texttt{tasks/audits/home\_price\_scale\_checks/code}. The three-page
\texttt{scale\_checks.pdf} includes raw gaps as well as these figures;
\texttt{scale\_checks.rds} and its standard data report retain the calculations.
Source snapshots are unchanged from the preceding packet; code is revision
\texttt{e76f4dd} plus the September 15 working-tree changes.
\normalsize
